% Section 4: Experimental Setup
\section{Experimental Setup}

\subsection{Research Questions}

We design experiments to answer two research questions:

\textbf{RQ1 (Mechanism):} Does error type determine localization reliability, and does unreliable localization cause gradient noise?

\textbf{RQ2 (Efficiency):} Does error-type gating improve training efficiency?

RQ1 validates the causal mechanism underlying our approach; RQ2 demonstrates the practical benefit.

\subsection{Hypothesis Structure}

We decompose RQ1 into four mechanism hypotheses:

\begin{itemize}
\item \textbf{H-M1:} Fine-grained penalties concentrate gradients at traceback locations
\item \textbf{H-M2:} Localization accuracy varies by error type (U\_line vs U\_ignore)
\item \textbf{H-M3:} Unreliable localization causes gradient concentration at wrong tokens
\item \textbf{H-M4:} Gating improves gradient signal-to-noise ratio
\end{itemize}

For RQ2, we test one existence hypothesis:

\begin{itemize}
\item \textbf{H-E1:} Error-type gating improves sample efficiency
\end{itemize}

Each hypothesis specifies success criteria and gate type (MUST\_WORK or SHOULD\_WORK). We require all MUST\_WORK hypotheses (H-M1, H-M3, H-E1) to pass for overall validation.

\subsection{Dataset}

We use the APPS dataset~\cite{hendrycks2021apps} for all experiments. APPS contains 5,000 training problems across three difficulty levels, providing sufficient volume for statistical power. The dataset's focus on algorithmic problems produces diverse error types, enabling meaningful comparison between U\_line and U\_ignore categories.

We sample 500 problems per error category for mechanism experiments (H-M1 through H-M4) and use the full training set for efficiency experiments (H-E1).

\subsection{Model}

We use CodeT5-small~\cite{wang2021codet5} (60M parameters) for proof-of-concept validation. While RLTF uses CodeT5-large (770M), the smaller model enables rapid iteration and demonstrates mechanism generality. We expect the gating mechanism to transfer to larger models---the underlying causal chain (traceback $\rightarrow$ penalty $\rightarrow$ gradient) is architecture-agnostic.

\textbf{Comparison approach note:} Our experiments validate the gating mechanism through controlled within-setup comparisons (fine-gated vs.\ fine-always vs.\ coarse-only) rather than direct comparison to RLTF's reported numbers. This design choice is deliberate: RLTF reports results using CodeT5-large (770M parameters) on APPS, while we use CodeT5-small (60M parameters) for proof-of-concept validation. The 13$\times$ parameter difference makes absolute accuracy comparisons uninformative---the models operate at fundamentally different capability levels. Instead, our experiments isolate the effect of error-type gating by comparing training conditions that differ only in feedback strategy, holding model architecture and scale constant.

\subsection{Metrics}

\textbf{Concentration ratio} (H-M1, H-M3): Gradient magnitude at error-line tokens divided by magnitude at non-error tokens. Values $>1$ indicate concentration at error location.

\textbf{Localization accuracy} (H-M2): Percentage of samples where traceback line matches ground-truth bug location (within $\pm 2$ lines tolerance).

\textbf{Signal-to-noise ratio} (H-M4): Gradient signal at ground-truth locations divided by signal at incorrect locations. Higher values indicate cleaner credit assignment.

\textbf{Steps to threshold} (H-E1): Training steps required to reach 30\% pass@1 on evaluation set. Lower values indicate better sample efficiency.

\subsection{Conditions}

We compare three feedback conditions:

\begin{itemize}
\item \textbf{Coarse-only:} Binary pass/fail reward, no token-level penalties
\item \textbf{Fine-always:} Standard RLTF with uniform fine-grained penalties
\item \textbf{Fine-gated:} Our approach---fine-grained for U\_line, coarse for U\_ignore
\end{itemize}

For mechanism experiments, we analyze gradient behavior under fine-always to characterize the noise problem, then compare to fine-gated to validate the solution.

\subsection{Ground Truth Annotation}

For H-M2 and H-M3, we require ground-truth bug locations to measure localization accuracy and gradient concentration at correct tokens. We use synthetic code templates with known bug locations, ensuring deterministic annotation. This simplifies validation but limits generalization claims---we note this as a limitation.
